\documentclass[11pt,a4paper]{article}
\usepackage[utf8]{inputenc}
\usepackage{amsmath}
\usepackage{geometry}
\geometry{a4paper, margin=1in}
\usepackage{graphicx}
\usepackage{hyperref}

\title{\textbf{Residuum\_Modulus: A Cyber-Physical System Security Agent for Networked Microgrids}}
\author{Gangaharidi Narayan Kashyap}
\date{\today}

\begin{document}

\maketitle

\section{System Overview}
The proposed architecture integrates a Large Language Model (Gemini 1.5 Pro) with a Physics-Informed Generative Adversarial Network (PI-GAN) to autonomously detect and mitigate Data Integrity Attacks in smart grid infrastructures. 

\section{Microgrid Topology}
The agent monitors a tripartite networked topology consisting of:
\begin{itemize}
    \item \textbf{Microgrid 1:} Residential (Photovoltaic \& Load dynamics)
    \item \textbf{Microgrid 2:} Commercial (Photovoltaic \& Load dynamics)
    \item \textbf{Microgrid 3:} Industrial (Photovoltaic \& Load dynamics)
\end{itemize}

\section{Operational Logic and MCP Integration}
The agent utilizes the MongoDB Model Context Protocol (MCP) Server to access historical load series via the \texttt{query\_historical\_load} tool. 

Upon retrieval, the agent employs a Chain-of-Thought reasoning framework to evaluate the data integrity. The \texttt{detect\_adversarial\_fault} tool triggers the PI-GAN inference model on Vertex AI. The PI-GAN evaluates the residual deviation between the incoming data $x$ and the learned physical manifold $\mathcal{M}$. 

If the confidence score of an adversarial perturbation exceeds the strict threshold of $\rho > 0.85$, the agent autonomously formulates a mitigation strategy, effectively bridging the gap between cyber-anomaly detection and physical grid resilience.

\end{document}